\begin{figure}[H]
\centering
\resizebox{\textwidth}{!}{%
\begin{tikzpicture}[
    node distance=2.4cm,
    % Data files: rounded gray boxes
    data/.style={
        draw, rounded corners=4pt, fill=gray!20,
        minimum width=2.0cm, minimum height=0.75cm,
        font=\small\ttfamily, align=center
    },
    % QC tools: blue boxes
    qc/.style={
        draw, fill=blue!20,
        minimum width=2.0cm, minimum height=0.75cm,
        font=\small, align=center
    },
    % Preprocessing tools: orange boxes
    preproc/.style={
        draw, fill=orange!30,
        minimum width=2.0cm, minimum height=0.75cm,
        font=\small, align=center
    },
    % Alignment: green boxes
    aligner/.style={
        draw, fill=green!25,
        minimum width=2.0cm, minimum height=0.75cm,
        font=\small, align=center
    },
    lbl/.style={font=\footnotesize\itshape, align=center, text width=2.2cm},
    arr/.style={-{Stealth[length=6pt]}, thick},
]

% ---- Nodes ----
\node[data]    (fastq)    {Raw\\FASTQ};
\node[qc]      (fqc1)     [right=of fastq]    {FastQC};
\node[preproc] (cutadapt) [right=of fqc1]     {Cutadapt};
\node[qc]      (fqc2)     [right=of cutadapt] {FastQC};
\node[aligner] (aligner)  [right=of fqc2]     {Aligner\\(BWA/STAR)};
\node[data]    (bam)      [right=of aligner]   {BAM};

% ---- Arrows ----
\draw[arr] (fastq)    -- (fqc1);
\draw[arr] (fqc1)     -- (cutadapt);
\draw[arr] (cutadapt) -- (fqc2);
\draw[arr] (fqc2)     -- (aligner);
\draw[arr] (aligner)  -- (bam);

% ---- Labels below each node ----
\node[lbl, below=0.25cm of fastq]    {input reads};
\node[lbl, below=0.25cm of fqc1]    {assess quality};
\node[lbl, below=0.25cm of cutadapt]{trim adaptors\\low-$Q$ bases};
\node[lbl, below=0.25cm of fqc2]    {confirm\\trimming};
\node[lbl, below=0.25cm of aligner] {map to\\reference};
\node[lbl, below=0.25cm of bam]     {aligned reads};

% ---- Legend ----
\begin{scope}[shift={(0, -2.2)}]
    \node[data,    minimum width=0.6cm, minimum height=0.45cm] (ld) at (0.3,   0) {};
    \node[anchor=west, font=\footnotesize] at (0.65, 0) {Data file};

    \node[qc,      minimum width=0.6cm, minimum height=0.45cm] (lq) at (3.2,   0) {};
    \node[anchor=west, font=\footnotesize] at (3.55, 0) {QC (inspect only)};

    \node[preproc, minimum width=0.6cm, minimum height=0.45cm] (lp) at (7.2,   0) {};
    \node[anchor=west, font=\footnotesize] at (7.55, 0) {Preprocessing (modifies reads)};

    \node[aligner, minimum width=0.6cm, minimum height=0.45cm] (la) at (12.4,  0) {};
    \node[anchor=west, font=\footnotesize] at (12.75, 0) {Alignment};
\end{scope}

\end{tikzpicture}%
}% end resizebox
\caption{Typical NGS preprocessing pipeline before alignment.
\textbf{FastQC} (blue) is a quality control tool --- it inspects and reports read
quality but does not modify any data.
\textbf{Cutadapt} (orange) is a \emph{preprocessing} tool --- its primary purpose
is to remove adaptor sequences and trim low-quality bases \emph{before alignment},
so that the aligner receives clean reads and produces accurate mappings.
A second FastQC run after trimming confirms that the cleaning was successful.
In this dataset the BAM file was provided pre-aligned, so Cutadapt was run
retrospectively on the raw FASTQs to confirm adaptor presence and measure
insert size, rather than as a mandatory preprocessing step before mapping.}
\label{fig:ngs-pipeline}
\end{figure}
